% \iffalse meta-comment
% Class for reports for the Computer Security course of the University of Namur
%
% Copyright (C) 2023 by Martin Jacob <martin.jacob@student.unamur.be>
% -------------------------------------------------------------------
% This work may be distributed and/or modified under the
% conditions of the LaTeX Project Public License, either version 1.3
% of this license or (at your option) any later version.
% The latest version of this license is in
%   http://www.latex-project.org/lppl.txt
% and version 1.3 or later is part of all distributions of LaTeX
% version 2005/12/01 or later.
%
% This work has the LPPL maintenance status `maintained'.
% 
% The Current Maintainer of this work is Martin Jacob.
%
% This work consists of the files secreport.dtx and secreport.ins
% and the derived file secreport.sty.
% \fi

% \iffalse
%<*driver>
\ProvidesFile{secreport.dtx}
%</driver>
%<class>\NeedsTeXFormat{LaTeX2e}
%<class>\ProvidesClass{secreport}
%<*class>
    [2023/02/15 v1.00 Reports for the Computer Security course of the University of Namur]
%</class>
%<*driver>
\documentclass[a4paper]{ltxdoc}

\usepackage[utf8]{inputenc}
\usepackage[T1]{fontenc}
\usepackage{hypdoc}
\usepackage[left=4cm,right=4cm]{geometry}
\usepackage[english]{babel}
\usepackage{verbatim}
\usepackage{csquotes}
\usepackage{enumitem}
\usepackage{xcolor}
\usepackage{etoolbox}
\usepackage{listofitems}
\usepackage{dirtree}

\EnableCrossrefs
\CodelineIndex
\RecordChanges

\hypersetup{
    colorlinks=true,
    urlcolor=cyan,
    pdfinfo={
        Title={secreport.cls},
        Subject={secreport class documentation file},
        Keywords={LaTeX, TeX, class, documentation, University of Namur, Computer Security, security},
        Author={Martin Jacob}
    }
}
\MakeOuterQuote{"}
\setlist[enumerate,itemize]{align=left}

\begin{document}
    \DocInput{\jobname.dtx}
    \newpage
    \PrintChanges
    \PrintIndex
\end{document}
%</driver>
% \fi
% 
% \GetFileInfo{\jobname.dtx}
%
% \def\DescribeEnvNoIndex{\leavevmode\begingroup\MakePrivateLetters
%  \Describe@EnvNoIndex}
% \def\Describe@EnvNoIndex#1{\endgroup
%               \marginpar{\raggedleft\PrintDescribeEnv{#1}}^^A
%               \ignorespaces}
%
% \begingroup
% \makeatletter
% \lccode`9=32\relax
% \lowercase{^^A
%   \edef\nosp{\noexpand\DoNotIndex{\@backslashchar9}}^^A
% }^^A
% \expandafter\endgroup\nosp
%
% \expandafter\DoNotIndex\expandafter{\bslash}
% \DoNotIndex{\@empty,\@gobbletwo,\\,\arabic,\begin,\bfseries,\ClassInfo,\csname,\CurrentOption,\DeclareRobustCommand,\def,\detokenize,\else,\end,\endcsname,\expandafter,\fi,\foreach,\gdef,\global,\huge,\href,\if,\IfEndWith,\ifnumcomp,\ifx,\ignorespaces,\ignorespacesafterend,\include,\includegraphics,\LARGE,\Large,\large,\let,\linewidth,\MakeUppercase,\mbox,\newcommand,\newenvironment,\newpage,\null,\numexpr,\par,\PassOptionsToClass,\ProcessOptions,\quad,\relax,\rule,\sectionignore,\textbf,\textcolor,\textit,\textsc,\textsuperscript,\texttt,\the,\thepage,\thesection,\thesubsection,\trim@spaces,\undefined,\vfill,\vskip,\vspace,\x,\xdef}
% 
% \title{\textsf{secreport} --- Reports for the Computer Security course of the University of Namur}
% \author{Martin Jacob \\ \href{mailto:martin.jacob@student.unamur.be}{\textcolor{black}{\texttt{martin.jacob@student.unamur.be}}}}
% \date{2023-02-15\\ v1.00}
% 
% \newcommand{\ctan}[1]{\href{https://ctan.org/pkg/#1}{\texttt{#1}}}
% \newcommand{\unamur}[2]{\href{#1}{\textcolor{unamur_color}{#2}}}
% \newcommand{\textopt}[1]{$\langle #1\rangle$}
% \newcommand*{\ms}[2]{\readlist*\msopt{#2}\relax\fbox{\cmd{#1\ifnumcomp{\msoptlen}{>}{0}{\foreachitem\x\in\msopt[]{\{\textopt{\x}\}}}{}}}\def\msopt{\undefined}\par}
% \newcommand{\es}[2]{\fbox{\parbox[t]{.4\linewidth}{\cmd{\begin\{#1\}\ifstrempty{#2}{}{\{$\langle #2\rangle$\}}\\...\\\bslash end\{#1\}}}}\par}
% \newenvironment{macrodesc}{\setlength{\parskip}{1ex plus .5ex minus .2ex}\setlength{\parindent}{0pt}\par\ignorespaces}{\par\ignorespacesafterend}
% \definecolor{unamur_color}{HTML}{69be28}
% 
% \maketitle
% \tableofcontents
% 
% \changes{v1.00}{2023/02/15}{First public release}
% 
% \section{Introduction}
% This \LaTeXe\ class was designed exclusively for the creation of reports for the practical classes of the Computer Security (INFOB301) course of the \unamur{https://unamur.be}{University of Namur}.
% It contains a predefined design, as well as some pre-configured packages considered useful for writing such reports (e.g., \ctan{minted}, \ctan{graphicx}, \ctan{hyperref}). It is based on the \ctan{article} class.
% 
% This class is somewhat opinionated regarding the structure of your report. It is strongly recommended that each challenge is in a separate file\footnote{A separate "challenges" folder can be set.}. You can find an example of this in the \autoref{subsubsec:struct-example}.
% 
% While this documentation is in English, the documents created with this class will be in French, for now. Support for English is planned in a future update.
% 
% \section{Usage}
% \subsection{Loading}\label{subsec:usage_loading}
% You can load this class the same way you load any document class:
% \\\fbox{\cmd{\documentclass[\textopt{options}]\{secreport\}}}
% 
% Regardless of the method used to get this class, you should have the \texttt{unamur.png}\footnote{SHA256: C8A7B146F2384DC4460558D1226DD816B97AE529B79BEFF36C8692E4D0234E46} file. If you do not have it for some reason, you can download it \href{https://pds.unamur.be/presse/logos/unamur.png}{here}.
% It is used in the title page and the header and is displayed with the \ctan{graphicx} package. It is thus required for compilation.
% 
% You can place it wherever you want, but if you decide to move it elsewhere than alongside your main document file, you have to set the designated folder using \ctan{graphicx}'s \cmd{\graphicspath}.
% 
% \subsection{Making the title page}
% \begin{macrodesc}
% \DescribeMacro{\and}{Even though the authors are now displayed vertically, the \cmd{\author} macro has not changed. Instead, \cmd{\and} now simply puts a classic line break (\cmd{\\}) for the tabular to typeset the authors as wanted.}
% 
% \DescribeMacro{\reportnum}
% {^^A
%     \ms{\reportnum}{num}
%     where \textopt{num} is the number of the current report.
    
%     Use this to set the report number that will be reflected in the title and the header. Defaults to 1.
% }
% 
% \DescribeMacro{\groupnum}
% {^^A
%     \ms{\reportnum}{num}
%     where \textopt{num} is the number of your group.
% 
%     Use this to set the group number that will be reflected in the header. Defaults to 1.
% }
% 
% \DescribeMacro{\maketitle}
% {^^A
%     Once \cmd{\reportnum} and \cmd{\groupnum} are set, use this to generate the title page. Just like in the \ctan{article} class, this has to be inside the \texttt{document} environment.
% }
% \end{macrodesc}
% 
% \subsection{Sectioning}
% \begin{macrodesc}
% \DescribeMacro{\section}
% {^^A
%     \ms{\section}{title}
%     where \textopt{title} is the title of your section.
% 
%     The sections have been thought to be exclusively used for sectioning challenges and thus, have their numbering removed.
% 
%     The main reason behind that reasoning is that you most likely want to write <<~Challenge 1~>> rather than <<~1.\quad Challenge 1~>>, but still having the section to be displayed in the table of contents, and still having the section numbering to be incremented for the subsections to be numbered correctly.
%  
%     As it is restricting, a separate command for this purpose while keeping the original sectioning format is not out of the question on a future update.
% }
% 
% \DescribeMacro{\tableofcontents}
% {^^A
%     Generates the table of contents on a separate page.
% }
% \end{macrodesc}
% 
% \subsection{Date}
% \begin{macrodesc}
% \DescribeMacro{\todayinquarters}
% {^^A
%     Prints the current date in quarters, according to the usual timings provided by the \unamur{https://www.unamur.be/etudes/etudiant/calendrier}{official academical calendar}. Might be wrong by a few days.
% 
%     \textit{e.g., <<~2\textsuperscript{e}\ quadrimestre 2021-2022~>>}
% }
% 
% \DescribeMacro{\date}
% {^^A
%     Same definition as the \ctan{article} class. Defaults to \cmd{\todayinquarters}
% }
% \end{macrodesc}
% 
% \subsection{Highlighting source code}
% As this work will be used in the context of a computer security course, you might want to include source code with syntax highlighting.
% 
% To do so, the \ctan{minted} package is imported with some default settings to help you typeset your code.
% What is provided here is a \textit{really} short documentation compared to the capacities of the package. It is advised to read the \href{http://mirrors.ctan.org/macros/latex/contrib/minted/minted.pdf}{original documentation} to master it properly.
% 
% \begin{macrodesc}
% \DescribeEnvNoIndex{minted}
% {^^A
%     \es{minted}{language}
%     where \textopt{language} is the language you want to set the syntax highlighting to.
% }
% 
% \DescribeMacro{\mintinline}
% {^^A
%     \ms{\mintinline}{language, code}
%     where \textopt{language} is the language you want to set the syntax highlighting to, and \textopt{code} is the line of code.
% 
%     Use this to typeset a single line of code inline. This box will not break, so be careful about overfull boxes.
% 
%     The given syntax is simplified for a common use case, as it is possible to set options and a custom delimiter for the \textopt{code}.
% }
% 
% \DescribeMacro{\inputminted}
% {^^A
%     \ms{\inputminted}{language, filename}
%     where \textopt{language} is the language you want to set the syntax highlighting to, and \textopt{filename} is the name of the file you want to import.
% 
%     The given syntax is simplified for a common use case, as it is possible to set options.
% }
% \end{macrodesc}
% 
% \subsection{Text formatting}
% This class sets \cmd{\parindent} to \fbox{2em} and \cmd{\parskip} to \fbox{1ex plus 0.5ex minus 0.2ex} by default.
% 
% To typeset paragraphs without a space between them, and without having to play with the \cmd{\parskip} length, a new environment is defined.
% 
% \begin{macrodesc}
% \DescribeEnv{noparskip}
% {^^A
%     \es{noparskip}{}
% }
% 
% \DescribeMacro{\code}
% {^^A
%   \ms{\code}{text}
%   where \textopt{text} is the text to typeset.
%   
%   Use this to quickly typeset simple code-like text (e.g., a filename with extension)
% }
% \end{macrodesc}
% 
% \subsection{Challenges}
% \begin{macrodesc}
% \DescribeMacro{\challengespath}
% {^^A
%   \ms{\challengespath}{path}
%   where \textopt{path} is the relative path to the folder containing the challenges files.
%   
%   Use this to configure the challenges folder for the \cmd{\includechallenges} command. Space are allowed inside \textopt{path}.
% }
% 
% \DescribeMacro{\challengespath}
% {^^A
%   \ms{\challengesprefix}{prefix}
%   where \textopt{prefix} is the prefix of your challenges files.
%   
%   Use this to configure the prefix of your files for the \cmd{\includechallenges} command. Space are allowed inside \textopt{prefix}.
% }
% 
% \DescribeMacro{\includechallenges}
% {^^A
%   \ms{\includechallenges}{num}
%   where \textopt{num} is the number of challenges to include.
% 
%   Use this to automatically include the challenges files either contained in the same folder as the file you're using this command in, or in the folder configured with \cmd{\challengespath}.
% }
% \end{macrodesc}
% 
% \subsubsection{Example}\label{subsubsec:struct-example}
% \begin{macrodesc}
% Let's say you decided to place your challenges files in a folder named "\texttt{Challenges}", each file beginning with "C", followed by the number of the challenge. You should then have a file tree structure similar to the following:
% 
% \fbox{\parbox[t]{.6\linewidth}{^^A
% \dirtree{^^A
%   .1 ./.
%   .2 Challenges.
%   .3 C1.tex.
%   .3 C2.tex.
%   .3 C3.tex.
%   .3 C4.tex.
%   .2 Images.
%   .3 unamur\_logo.png\DTcomment{Refer to \nameref{subsec:usage_loading}}.
%   .2 main.tex\DTcomment{Your main document file}.
% }}}
% 
% Then you would have to type:
% \end{macrodesc}
% \begin{itemize}
%   \item \cmd{\challengespath\{Challenges\}}
%   \item \cmd{\challengesprefix\{C\}}
% \end{itemize}
% in the preamble and
% \begin{itemize}
%   \item \cmd{\includechallenges\{4\}}
% \end{itemize}
% after \cs{begin\{document\}}.
% 
% \section{Class options}
% \begin{enumerate}
%     \item[\texttt{nogeometry}] prevents the loading of the \ctan{geometry} package. This exists as long as a way for users to load with custom options the \ctan{geometry} package on their own has not been implemented. For now, this class loads it with options "a4paper,top=25mm" by default.
% \end{enumerate}
% The remaining options are the same as the \ctan{article} class.
% 
%\StopEventually{^^A
%  \PrintChanges
%  \PrintIndex
%}
% 
% \section{Implementation}
% \subsection{Options, Base class and Required packages}
%    \begin{macrocode}
%<*class>
\DeclareOption{nogeometry}{\def\secreport@opt@nogeo}
\DeclareOption*{\PassOptionsToClass{\CurrentOption}{article}}
\ProcessOptions\relax
\LoadClass[a4paper]{article}

\RequirePackage{hyperref}
\RequirePackage{fancyhdr}
\RequirePackage{graphicx}
\RequirePackage{minted}
\RequirePackage{etoolbox}
\RequirePackage{xstring}
\RequirePackage{pgffor}
\RequirePackage{trimspaces}
\ifx\secreport@opt@nogeo\undefined
    \RequirePackage[a4paper,top=25mm]{geometry}
    \ClassInfo{secreport}{
        The 'nogeometry' option has not been used.
        The geometry package has been loaded with options "a4paper,top=25mm".}
\fi
\let\secreport@opt@nogeo\@empty
%    \end{macrocode}
% 
% \subsection{Title}
% \iffalse

%% === Title ===
% \fi
% \begin{macro}{\reportnum}
% Define a setter for the internal macro for the report number, used in the title and the header. Takes a number as argument.
%    \begin{macrocode}
\DeclareRobustCommand*{\reportnum}[1]{\gdef\@reportnum{#1}}
%    \end{macrocode}
% \end{macro}
% \begin{macro}{\groupnum}
% Define a setter for the internal macro for the group number, used in the header. Takes a number as argument.
%    \begin{macrocode}
\DeclareRobustCommand*{\groupnum}[1]{\gdef\@groupnum{#1}}
%    \end{macrocode}
% \end{macro}
% \begin{macro}{\and}
% Redefines |\and| to allow for vertical author list in the title.
%    \begin{macrocode}
\gdef\and{\\}
%    \end{macrocode}
% \end{macro}
% \iffalse

% \fi
% \begin{macro}{\@getquarters}
% The |\todayinquarters| command will need to get the current quarter.
%    \begin{macrocode}
\newcommand*{\@getquarters}
{%
    \ifnumcomp{\the\month}{>}{1}
    {%
        \ifnumcomp{\the\month}{>}{6}
        {%
            \ifnumcomp{\the\month}{>}{8}
            {%
                1%
            }{3}%
        }{2}%
    }{1}%
}
%    \end{macrocode}
% \end{macro}
% \begin{macro}{\@getschoolyear}
% The |\todayinquarters| command will need to get the current school year.
%    \begin{macrocode}
\newcommand*{\@getschoolyear}
{%
    \ifnumcomp{\the\month}{>}{8}
    {%
        \ifnumcomp{\the\day}{>}{13}
        {%
            \the\year-\the\numexpr(\the\year+1)%
        }{\the\numexpr(\the\year-1)\relax-\the\year}%
    }{\the\numexpr(\the\year-1)\relax-\the\year}%
}
%    \end{macrocode}
% \end{macro}
% \begin{macro}{\todayinquarters}
% Define a command to typeset the current date as quarters.
%    \begin{macrocode}
\newcommand*{\todayinquarters}
{%
    \@getquarters\ifnumcomp{\@getquarters}{=}{1}
        {\textsuperscript{er}}
        {\textsuperscript{e}}
    quadrimestre \@getschoolyear%
}
%    \end{macrocode}
% \end{macro}
% 
% \iffalse

% \fi
% Set default values.
%    \begin{macrocode}
\title{Rapport du TP n°}
\reportnum{1}
\groupnum{1}
\date{\todayinquarters}
%    \end{macrocode}
% 
% \iffalse

% \fi
% \begin{macro}{\maketitle}
% Redefine |\maketitle| to generates the design and typeset accordingly on a dedicated page.
%    \begin{macrocode}
\renewcommand{\maketitle}
{%
    \begin{noparskip}
        \begin{titlepage}%
            \begin{center}%
                \textsc{\LARGE\href{https://unamur.be}{\textcolor{black}{Université de Namur}}}%
                \vspace{1.5cm}%
                
                \textsc{\Large Sécurité Informatique}%
                \vspace{1cm}%
                
                % Title
                \HRule \vspace{.5cm}%
                \textbf{\huge \@title\@reportnum}\\%
                \HRule \vspace{1cm}%
                
                % UNamur Logo
                \includegraphics[height=180px]{unamur}%
                \vspace{1cm}%
                
                % Date
                {\LARGE \@date}%
                \vspace{4cm}%
                
                % Group members (Author(s))
                \underline{\textbf{\textsc{Groupe \@groupnum}}}%
                \vspace{.5em}%
                
                \begin{tabular}[t]{c}%
                    \@author
                \end{tabular}%
            \end{center}%
            \vfill\null
        \end{titlepage}%
    \end{noparskip}
    \global\let\maketitle\relax
    \global\let\@title\@empty
    \global\let\@author\@empty
    \global\let\@date\@empty
    \global\let\@thanks\@empty
    \global\let\title\relax
    \global\let\author\relax
    \global\let\date\relax
    \global\let\thanks\relax
    \global\let\and\relax
}
%    \end{macrocode}
% \end{macro}
% 
% \subsection{Table of Contents}
% \iffalse

%% === Table of Contents ===
% \fi
% \begin{macro}{\tableofcontents}
% Redefine |\tableofcontents| to generate the table of contents on a dedicated page with no space between items.
%    \begin{macrocode}
\let\latex@tableofcontents\tableofcontents
\renewcommand{\tableofcontents}
{%
    \newpage%
    \begin{noparskip}%
        \latex@tableofcontents%
    \end{noparskip}%
    \vfill%
}
%    \end{macrocode}
% \end{macro}
% 
% \subsection{Page Style}
% \iffalse

%% === Page Style ===
% \fi
% Set the page style to \texttt{fancy} for the \ctan{fancyhdr} package commands.
%    \begin{macrocode}
\pagestyle{fancy}
%    \end{macrocode}
% \subsubsection{Header/Footer}
% \iffalse

%% == Header/Footer ==
% \fi
% Clear the header and the footer.
%    \begin{macrocode}
\fancyhf{}
%    \end{macrocode}
% Configure the header.
%    \begin{macrocode}
\fancyhead[L]{\includegraphics[trim={0 200px 0 0},clip,height=30px]{unamur}}
\fancyhead[C]{\large{\textbf{Rapport du TP n°\@reportnum\ - Groupe \@groupnum}}}
\fancyhead[R]{INFOB301\\\href{https://unamur.be}{\textcolor{black}{UNamur}}}
\setlength{\headheight}{40pt}
%    \end{macrocode}
% \iffalse

% \fi
%% Remove extra space in the section name display.
%    \begin{macrocode}
\renewcommand*{\sectionmark}[1]{\markboth{\MakeUppercase{#1}}{}}
%    \end{macrocode}
% Configure the footer.
%    \begin{macrocode}
\fancyfoot[L]{\textcolor{gray}{\leftmark}}
\fancyfoot[R]{\textcolor{gray}{\thepage}}
\renewcommand*{\footrulewidth}{0.5pt}
%    \end{macrocode}
% 
% \subsection{Sections}
% \iffalse

%% === Sections ===
% \fi
% Removing section numbering while keeping the counter for the subsections to display properly.
%    \begin{macrocode}
\renewcommand*{\thesection}{}
\renewcommand*{\thesubsection}{\arabic{section}.\arabic{subsection}}
\let\sectionignore\@gobbletwo
\def\@seccntformat#1{\csname #1ignore\expandafter\endcsname\csname the#1\endcsname\quad}
\let\latex@numberline\numberline
\def\numberline#1{\if\relax#1\relax\else\latex@numberline{#1}\fi}
%    \end{macrocode}
% 
% \subsection{Text Formatting}
% \iffalse

%% === Text Formatting ===
% \fi
% \begin{macro}{\code}
% Define a macro to quickly typeset simple code-like text.
%    \begin{macrocode}
\newcommand*{\code}[1]{\mbox{\texttt{\textit{#1}}}}
%    \end{macrocode}
% \end{macro}
% 
% \iffalse

% \fi
% \begin{macro}{\HRule}
% Define a macro to make a full line thin rule. Used in the title generation.
%    \begin{macrocode}
\newcommand*{\HRule}{\rule{\linewidth}{0.5mm}\par}
%    \end{macrocode}
% \end{macro}
% 
% \iffalse

% \fi
% Set default values for paragraph indent and in-between paragraph spacing.
%    \begin{macrocode}
\setlength{\parindent}{2em}
\setlength{\parskip}{1ex plus 0.5ex minus 0.2ex}
%    \end{macrocode}
% 
% \iffalse

% \fi
% Set default configuration for the \ctan{hyperref} package.
%    \begin{macrocode}
\hypersetup{
    colorlinks=true,
    linkcolor=blue,
    urlcolor=blue
}
%    \end{macrocode}
% 
% \iffalse

% \fi
% Set default configuration for the \ctan{minted} package.
%    \begin{macrocode}
\setminted{autogobble, bgcolor=lightgray, breaklines}
\setminted[Python]{linenos, frame=lines, framesep=2mm}
%    \end{macrocode}
% \iffalse

% \fi
% \begin{macro}{noparskip}
% Define an environment to typeset paragraphs without in-between paragraph spacing.
%    \begin{macrocode}
\newenvironment{noparskip}
{%
    \vspace{\parskip}%
    \setlength{\parskip}{0pt}%
    \par\ignorespaces%
}
{\par\ignorespacesafterend}
%    \end{macrocode}
% \end{macro}
% 
% \subsection{Challenges}
% \iffalse

%% === Challenges ===
% \fi
% ^^A \newcommand{\challenge}[1]{\include{#1}}
% \begin{macro}{\challengespath}
% Define a command to set the path to the folder containing the challenges files.
%    \begin{macrocode}
\DeclareRobustCommand*{\challengespath}[1]{\gdef\@challengespath{\trim@spaces{#1}}}
%    \end{macrocode}
% \end{macro}
% \begin{macro}{\challengesprefix}
% Define a command to set the prefix of the challenges files.
%    \begin{macrocode}
\DeclareRobustCommand*{\challengesprefix}[1]{\gdef\@challengesprefix{#1}}
%    \end{macrocode}
% \end{macro}
% \iffalse

% \fi
% \begin{macro}{\includechallenges}
% Define a command to include a specified number of challenges files. Takes a number as argument. The path to the files can be set with the |\challengespath| command.
%    \begin{macrocode}
\newcommand{\includechallenges}[1]
{%
    \foreach \x in {1,...,#1}{%
        \IfEndWith{\@challengespath}{/}
            {\include{\@challengespath \@challengesprefix\x}}
            {\include{\@challengespath/\@challengesprefix\x}}%
    }%
    \global\let\challengespath\relax%
    \global\let\@challengespath\@empty%
}
%</class>
%    \end{macrocode}
% \end{macro}
% \newpage
% \Finale